\section{Discussion}\label{sec:discussion}

\emph{When the layer acts.} The effect depends on two properties of the
pipeline, not on the dataset: whether the tracker has a continuation step,
and whether the detector's scores match the tracker's thresholds. On MOT17
with well-matched detections, the two-stage trackers BoostTrack, ByteTrack at
floor 0.1, and Hybrid-SORT produced output identical to their hosts, while
the single-stage OC-SORT and PD-SORT gained HOTA on every sequence. On the
external streams declared after the freeze, the layer classified at least
99.7\% of the frames as clean on MOT17 and DanceTrack, and it left those
pipelines at or near their own output. When the stream fits the tracker, the
layer mostly stays out of the way, and when the tracker lacks a step, the
layer adds it.

\emph{Self-limiting by design.} In the clean mode the layer keeps the
tracker's thresholds and scores, so a tracker with a low-score step and
well-matched detections receives its normal input. This is the main
difference from the earlier design \aob{}, which lowered the HOTA of the same
strong trackers by 4.15 to 5.89 points (Table~\ref{tab:prior}). The layer
therefore does not need to know in advance whether a tracker needs help; it
reads the answer from the scores.

\emph{Detector replacement.} When a detector is replaced or its score scale
shifts, the fixed thresholds of a tracker no longer match the scores. The
layer adjusts them from the scores without labels. Under halved scores it
restored two trackers from 0 to about 66 HOTA, and with a different detector
it raised the MOTA of two two-stage trackers by 6.0 and 7.9 points. This is
the case the layer was built for: a detector that is swapped, retrained, or
quantized while the tracker stays the same.

\emph{Failure cases.} Three results go against the layer. First, with the
under-confident YOLOv8n detector on KITTI, the noisy mode restricted new
tracks of the two two-stage trackers and cost them 1.87 and 2.43 MOTA points
(Sect.~\ref{sec:detector}). Second, the loss of C-TWiX on KITTI cars
($-1.52$ HOTA) comes almost entirely from one short sequence: in sequence 0014
the layer classified 45 of 106 frames as noisy, and HOTA fell by 17.09 points,
while six of the other eight sequences changed by less than 0.4 points. The
same sequence causes the loss on KITTI pedestrians (8.50 points in that
sequence). In such a short sequence, the median of $\rho$ over the stream has
little history, and the restrictions of the noisy mode most likely removed
detections that the coordinate-based association of C-TWiX relies on; this
mechanism was not isolated further. A rule that requires a minimum history
before a noisy decision is a natural remedy, but it was not part of the frozen
configuration. Third, C-TWiX on DanceTrack lost 1.00 HOTA with an interval
including zero; two sequences lost more than 10 points and one gained 9.37.
The stream was clean on almost every frame, so the effect comes from the
operations of the clean mode (cross-class duplicate removal, threshold
lowering to $\sigma(t_2)$, and the continuation rescue) acting on a tracker
that learns association from coordinates; which of them causes the loss was
not isolated.

\emph{Limitations.} The strongest evidence is development evidence: the
configuration was selected on MOT17, KITTI, and VisDrone, so the gains there
are not held out. Only four external systems passed their reproduction, and
only PD-SORT improved. The external tests change the tracker, and with C-TWiX
and TrackTrack the detector and the dataset, but they use the detections
released by each tracker's authors on public benchmarks, not new data. The
paired intervals are computed for each comparison separately, without a
correction for multiple comparisons, and with 7 to 25 sequences per cell they
are wide. Finally, the layer was compared with each tracker alone and with
its own earlier design, not with other adaptive methods such as Adaptive
ByteTrack \cite{ma2023adaptivebytetrack} or ACT \cite{act2026adaptive}, nor
with a threshold-only or a rescue-only variant of itself; these comparisons
are left for future work.

\emph{Relation to detector-side control.} Our earlier AC-MOT controller,
described in a manuscript submitted for publication, acts before the detector
and chooses its operating point from the scene. \scl{} acts after the
detector and adapts the tracker's input to the scores that the detector
produces. The two layers work on different parts of the pipeline, so they
could be used together; combining them is left for future work.

\emph{Cost.} The layer needs no training, no GPU, and no access to the inside
of the detector, and it needs only four contract values from the tracker's
configuration. At run time it reads the candidate list, one small grayscale
image for the motion cue, and the tracks of the previous frame. On a CPU it
takes 2.95 ms per frame, adds 4.5\% to the total time of the tested pipeline,
and makes the tracker itself faster. New trackers are added through their
published contract values and an adapter, without changes to the rules.
